%% Paper F: do rankings of explanation methods hold across tabular datasets?
%% Target: Journal of Computer Sciences Institute (JCSI). Rules: VENUE_JCSI.md.
%%
%% THIS IS THE SOURCE. paper_f.tex is GENERATED from it by scripts/build_paper_f.py, which
%% replaces \input{generated/...} fragments and <<alias:key:column:format>> placeholders
%% with values read from result files. Edit this file, never paper_f.tex.
%%
%% The journal takes Word in its own template: A4, two columns of 8 cm, 1 cm between,
%% margins 2.5 cm (top, bottom) and 2 cm (left, right), Times New Roman 10 pt, single
%% spacing, 4 to 8 pages. This source is set to the same page so that the PDF predicts
%% the length. The submitted file is made in the journal's Word template.
%%
%% Review is double-blind. The default build omits author identity; the full build is
%% made by defining \fullversion before this file.
%%
%% DRAFT of 2026-10-05: introduction and methods follow ANALYSIS_PLAN.md (version 2).
%% The experiment has not been run. Every result is a visible "[pending]" mark; no
%% number of the results, discussion or conclusions exists yet.
\documentclass[10pt,a4paper,twocolumn]{article}
\usepackage[top=25mm,bottom=25mm,left=20mm,right=20mm,columnsep=10mm]{geometry}
\usepackage{fontspec}
\setmainfont{texgyretermes}[Extension=.otf, UprightFont=*-regular, BoldFont=*-bold,
  ItalicFont=*-italic, BoldItalicFont=*-bolditalic]
\usepackage{amsmath}
\usepackage{graphicx}
\usepackage{booktabs}
\usepackage{array}
\usepackage{xcolor}
\usepackage[hidelinks]{hyperref}
\usepackage[numbers,sort&compress]{natbib}
\setlength{\emergencystretch}{3em}
\setlength{\bibsep}{2pt}
\providecommand{\fullurl}[1]{\url{#1}}  % used by entries taken from the Paper E list
\newcommand{\pending}[1]{\textcolor{red}{[pending: #1]}}
%% Code and data archive (full version only). Set when the Zenodo version exists.
\newcommand{\paperfarchive}{\pending{Zenodo DOI}}

\newif\iffull
\ifdefined\fullversion \fulltrue \else \fullfalse \fi

\begin{document}

\twocolumn[
\begin{flushleft}
{\Large\bfseries Comparative analysis of four explanation methods across
<<sum:datasets:value:w>> tabular datasets: do benchmark rankings generalise?\par}
\medskip
\iffull
Jonathan Herrera-V\'asquez\textsuperscript{*}\par
\smallskip
{\small Department of Computer Science, Universidad Americana de Europa (UNADE),
Canc\'un, Quintana Roo, Mexico\par}
\else
\textit{Author information omitted for double-blind review.}\par
\fi
\medskip
\textbf{Abstract}\par
\smallskip
Explanation methods are usually compared on one dataset, and the resulting order is read
as a property of the methods. We test whether it is. LIME, SHAP, Anchors and DiCE explain
four classifiers on <<sum:datasets:value:d>> public tabular datasets chosen by a fixed
rule; each method is ranked on faithfulness, stability, sparsity and cost. The agreement
of the rankings between datasets is the mean rank correlation, with a bootstrap interval
and a permutation test. \pending{main result, one sentence per ranked measure}
\pending{what a result on the Adult dataset does and does not predict}

\medskip
\textit{Keywords}: explainable artificial intelligence; LIME; SHAP; benchmark;
external validity; tabular data
\medskip

\iffull
{\small \textsuperscript{*}Corresponding author. \textit{Email address}:
jonnabio@gmail.com (J. Herrera-V\'asquez). ORCID:
\url{https://orcid.org/0000-0002-7149-6635}\par}
\medskip
\fi
\hrule
\bigskip
\end{flushleft}
]

\section{Introduction}

A model that decides on credit, hiring or medical risk is often too complex to be read
directly, and post-hoc explanation methods are used to say which inputs drove a
prediction. LIME fits a local linear model around the instance~\cite{ribeiro2016why},
SHAP distributes the prediction among the features with Shapley
values~\cite{lundberg2017unified}, Anchors returns a rule that keeps the
prediction with high probability~\cite{ribeiro2018anchors}, and DiCE returns
counterfactual instances that change it~\cite{mothilal2020dice}. Surveys of the field
describe these methods and list the evaluation of explanations among its open
problems~\cite{longo2024manifesto,mersha2024survey}, also for tabular
data~\cite{obrienquinn2024tabular,santana2026tabular}. A practitioner has to choose
among the methods, and comparative studies are the usual guide.

\subsection{Related work}

There is no single accepted measure of the quality of an explanation. A study of the
evaluation of local post-hoc methods on tabular data derives 20 criteria with guidelines
for each~\cite{velmurugan2025guidelines}; an examination of the many proposed metrics
finds duplication and calls for standard frameworks~\cite{pawlicki2024metrics}; and new
measures, benchmarks and metric studies continue to
appear~\cite{bello2025strength,sithakoul2024beexai,pereira2026local}. Where a ground
truth for the explanation can be built, methods differ clearly in
fidelity~\cite{mironicolau2024fidelity}.

The methods compared here have known weaknesses. SHAP and LIME depend strongly on the
model and on collinear features~\cite{salih2024perspective}, they can assign importance
to features that a model with known rules does not use~\cite{hooshyar2024problems}, and
Shapley values have been criticised as a basis for
explanations~\cite{huang2024failings}. Different methods can explain the same
prediction differently, and measures of their agreement have been
proposed~\cite{parashar2026agreement,gao2026uncertainty}. Rule-based and counterfactual
explanations for tabular data are active lines of
work~\cite{tekin2026fuzzy,redelmeier2024mcce}. Applied studies compare the methods on
one task: three methods on five datasets of different types in this
journal~\cite{kuszewska2026comparative}, four importance measures on one medical
dataset~\cite{biesaga2026comparison}, and LIME and SHAP on the prediction of software
issues~\cite{schulte2024studying}.

A comparison made on one dataset leaves a question open: would the order of the methods
be the same on another? For attribution maps of images the answer was negative: metric
results on one dataset did not necessarily hold on
others~\cite{gevaert2024evaluating}. A recent benchmark of tabular classifiers
curates a representative collection of datasets~\cite{erickson2025tabarena}, and
tabular data are heterogeneous in feature types, cardinality and
sparsity~\cite{somvanshi2026survey}. Comparisons of explanation methods on tabular data
seldom use many datasets. The census income data known as Adult is a frequent
single benchmark, and a framework for measuring explanation quality was itself first
validated on it~\cite{herrera2026framework}.

\subsection{Purpose and hypothesis}

The purpose of this work is to measure how far the order of four explanation methods,
obtained on one tabular dataset, holds on other tabular datasets. The object of the
research is therefore the transfer of a comparative conclusion and not the quality of
any one method. The research covers binary classification on tabular data, four
classifiers of different families and four measures on which the four methods can all be
ranked.

The research hypothesis is: \textit{on each measure, the order of the four explanation
methods agrees across tabular datasets more than chance.} The hypothesis, the measures,
the rule for choosing the datasets and the decision rule were fixed in a plan committed
to a public repository before any explanation was computed.

\section{Materials and methods}

\subsection{Datasets}

The datasets were chosen by a rule and not by hand, so that the choice could not favour
one outcome. A dataset is included if it is public and usable without registration; is
tabular with a binary target; has at least 1,000 rows and between 5 and 100 features; and
if at least three of the four classifiers reach a cross-validated area under the ROC
curve between 0.65 and 0.98 on it, which removes datasets that are noise or trivially
separable. Candidates were taken from OpenML~\cite{vanschoren2014openml} in an order
fixed beforehand, in six strata defined by the type of the features (numeric, mixed,
categorical) and the number of rows (below or from 10,000).
When fewer than sixteen of the listed candidates passed, the list was extended with the
binary tasks of the OpenML-CC18 suite in the order of their identifier. Of
<<sum:candidates:value:d>> candidates examined, <<sum:datasets:value:d>> are in the
study (Table~\ref{tab:datasets}). The minority class is the positive class.

\begin{table*}[t]
\centering
\caption{Datasets of the study. Features are counted before encoding.}
\label{tab:datasets}
\small
\begin{tabular}{llrrrrl}
\toprule
Dataset & OpenML id & Rows & Features & Numeric (\%) & Minority class (\%) & Stratum \\
\midrule
phoneme & 1489 & 5,404 & 5 & 100 & 29.3 & numeric, under 10,000 rows \\
qsar-biodeg & 1494 & 1,055 & 41 & 100 & 33.7 & numeric, under 10,000 rows \\
ozone-level-8hr & 1487 & 2,534 & 72 & 100 & 6.3 & numeric, under 10,000 rows \\
MagicTelescope & 1120 & 19,020 & 10 & 100 & 35.2 & numeric, from 10,000 rows \\
electricity & 151 & 45,312 & 8 & 88 & 42.4 & mixed, from 10,000 rows \\
default-of-credit-card-clients & 42477 & 30,000 & 23 & 100 & 22.1 & numeric, from 10,000 rows \\
credit-g & 31 & 1,000 & 20 & 35 & 30.0 & mixed, under 10,000 rows \\
churn & 40701 & 5,000 & 20 & 80 & 14.1 & mixed, under 10,000 rows \\
compas-two-years & 42192 & 5,278 & 13 & 54 & 47.0 & mixed, under 10,000 rows \\
adult & 1590 & 48,842 & 14 & 43 & 23.9 & mixed, from 10,000 rows \\
bank-marketing & 1461 & 45,211 & 16 & 44 & 11.7 & mixed, from 10,000 rows \\
Amazon\_employee\_access & 4135 & 32,769 & 9 & 0 & 5.8 & categorical, from 10,000 rows \\
pc4 & 1049 & 1,458 & 37 & 100 & 12.2 & numeric, under 10,000 rows \\
pc3 & 1050 & 1,563 & 37 & 100 & 10.2 & numeric, under 10,000 rows \\
jm1 & 1053 & 10,885 & 21 & 100 & 19.4 & numeric, from 10,000 rows \\
kc1 & 1067 & 2,109 & 21 & 100 & 15.5 & numeric, under 10,000 rows \\
\bottomrule
\end{tabular}
\end{table*}


Numeric features are standardised and missing values replaced by the median; categorical
features are one-hot encoded, with categories beyond the 20 most frequent of a feature
grouped into one level. Each dataset is split into 80\% for training and 20\% for
testing, stratified by class.

\subsection{Classifiers and explanation methods}

Four classifiers are trained on every dataset with one configuration for all: logistic
regression, a random forest of 100 trees,
XGBoost with 100 trees of depth 4, and a multilayer perceptron
with one hidden layer of 100 units.

The four explanation methods are also used with one configuration for all datasets and
no tuning: LIME with 1,000 samples, 10 features and a kernel width of 3; SHAP with 50
background rows, in its exact tree variant for the two tree models and its kernel variant
for the other two; Anchors with a precision threshold of 0.95; and DiCE with random
search for one counterfactual. Explanations are computed on the encoded features. An
Anchors rule is read as the set of features it uses, and a counterfactual as the features
it changes with the size of each change.

\subsection{Research scenario}

Each of the <<sum:datasets:value:d>> datasets is combined with the four classifiers and
the four methods, and every combination is repeated with five random seeds that fix the
split, the classifier and the sampling of the method. For each dataset, classifier and
seed, 200 test instances are drawn, stratified by predicted class, and the four methods
explain the same instances. An instance for which a method returns no explanation is kept
and counted as a failure.

\subsection{Measures}

Four measures are defined for all four methods and are used to rank them.
\textit{Faithfulness gap}: the absolute change of the predicted probability when the
features named by the explanation, up to the top 20\% of all features, are replaced by
their training mean or mode; higher is better. \textit{Stability}: the mean Jaccard
index between the five top features of the explanation and
those of the explanations of 10 perturbed copies of the instance, where numeric features
receive Gaussian noise of 0.05 standard deviations; higher is better. \textit{Sparsity}:
the share of features the explanation does not use; higher is better. \textit{Cost}: the
time to produce one explanation; lower is better. A failed instance takes the worst value
observed for its dataset and classifier, so that failures cannot improve a mean.

\subsection{Statistical analysis}

The unit of analysis is the dataset. For each dataset and measure, the values are
averaged over instances, seeds and classifiers, and the four methods are ranked. The
agreement between datasets is the mean Spearman correlation between the rankings of all
pairs of datasets, which is Kendall's coefficient of concordance corrected for the number
of datasets. It is reported with a 95\% percentile bootstrap
interval over datasets (10,000 resamples) and tested against
unrelated rankings with a permutation test (100,000 permutations of the ranks inside each
dataset), with the Holm correction over the four measures. The same
correlation between the five seeds inside a dataset gives a ceiling: agreement between
datasets cannot be expected to exceed agreement between two runs on the same dataset.

The decision rule was fixed in advance. With a corrected $p$ below 0.05, the order is
said to generalise if the interval lies above 0.5, to be shared only weakly if it lies
below 0.5, and to be of undetermined strength otherwise. A simulation made before the
experiment showed that sixteen datasets detect a true correlation of 0.5 in nearly every
study, whereas four detect it in half of them.

Three further analyses are reported: for each pair of methods, the share of datasets in
which the pair is ordered as on Adult, with an exact binomial interval; the largest rank
change of each method between two datasets; and the share of the variance of the measure
that belongs to the method, to the dataset and to their interaction, with the
interaction tested by permutation.

\section{Results}

\pending{The experiment has not been run. This section will hold: Table 2, mean rank of
each method per measure with the mean rank correlation, its interval, the corrected p,
the ceiling and the decision; Figure 1, the rank of each method on each dataset, one
panel per measure; Table 3, agreement with the order found on Adult for each pair of
methods; one paragraph on failures; one paragraph on the variance shares.}

\section{Discussion and conclusions}

\pending{Discussion of the hypothesis for each measure, against the ceiling.}

\textit{Possible sources of error.} \pending{to be completed with the results. Known
before the run: explanations are computed on encoded features, so a categorical variable
counts as several features; Anchors and DiCE are compared with LIME and SHAP through the
features they name, not through their native output; each method has one configuration, and with it the kernel variant of SHAP returns at
most 10 non-zero attributions while the tree variant returns one for every feature;
four of the datasets are software-defect data of one family, so the sixteen are less
independent than their number suggests;
the datasets are binary classification tasks from one repository; cost depends on the
machine.}

\pending{Three numbered conclusions and directions for further research.}

\section*{Data and code availability}
\iffull
All datasets are public and are identified by their OpenML identifier in the archive.
The plan, the code, the per-instance results and the scripts that produce every table
and figure are available at \url{https://github.com/jonnabio/xai-eval-framework},
archived at \paperfarchive.
\else
All datasets are public. The plan, the code, the per-instance results and the scripts
that produce every table and figure are in a public archive whose address is omitted for
double-blind review.
\fi

%% The journal requires this declaration, with the name of the tool and the reason for
%% its use. Text approved by the author "for now" on 2026-10-05; the author confirms it
%% before submission, and it is not changed without the author's approval.
\section*{Declaration on the use of generative AI}
The author used the AI tool Claude (Anthropic), under the author's direction, to help
write the analysis code, edit the English text and format the references. The research
question, the plan, the design and the interpretation are the author's. The author
reviewed all content and takes full responsibility for it.

%% Author's rule (2026-10-05) and the journal's template: every reference prints its DOI
%% as an active URL. The style writes \href{https://doi.org/...}{doi:...}; print the URL.
{\renewcommand{\href}[2]{\url{#1}}
\bibliographystyle{elsarticle-num}
\bibliography{references}}

\end{document}
